\section*{Discussion}

\mname\ connects bidirectional spectral translation with molecular representation learning. On the reported computed-spectrum benchmarks, the model improves reconstruction over the evaluated VAE and Transformer baselines. Transfer experiments show lower MAE across the seven held-out ViBench domains, although improvements in $R^2$ are not universal. Frozen embeddings also support molecular-descriptor prediction beyond the original translation task. Together, these observations suggest that learning IR--Raman correspondence can provide both estimates of a missing spectral modality and features useful for downstream analysis.

The contribution of flow matching requires a narrower interpretation than the aggregate benchmark comparison alone permits. SpectraDiT-Direct already achieves strong endpoint accuracy with one network evaluation. The eight-step Flow model performs better in the reported QM9S comparison, but its longer training budget and larger endpoint-loss coefficient prevent attribution of the difference solely to integration. Similarly, the VAE and Transformer comparisons change several architectural and optimization choices at once. The present results support the complete method under the stated protocols; they do not establish that flow integration is necessary for accurate translation or identify the isolated contribution of every loss component.

The learned trajectory offers a way to inspect the mapping between modalities. Intermediate states reveal where intensity changes as the model approaches its prediction, and peak-resolved analyses show that the largest residuals occur near spectral bands. These paths describe numerical transport through spectral space, not vibrational dynamics, reaction pathways or mode-resolved physical mechanisms. Their practical value is diagnostic: they expose the evolution of a prediction and help locate regions where peak position, width or intensity remains inaccurate. Raman$\rightarrow$IR is harder on the three principal benchmarks, but the transfer results show that this directional asymmetry depends on the dataset and metric.

The descriptor probes provide a separate test of representation utility. A frozen nonlinear encoder makes several structure-derived descriptors more accessible to a linear predictor than the raw spectral inputs do. This does not imply that the embedding creates molecular information absent from its source or that it replaces explicit structure. Morgan fingerprints remain a strong reference for descriptors computed from molecular graphs, and combining fingerprints with Flow features helps some tasks while degrading others. Moreover, the probes are fitted separately within each evaluation dataset. The results therefore demonstrate transfer of an encoder, rather than a property predictor that generalizes to a new domain without labelled examples.

Experimental evaluation places further limits on the current model. NIST IR inputs are compared with computed Raman targets, whereas the OpenSpecy/RRUFF study uses a separately trained model and identity-matched spectra from different library records. Neither is a same-specimen, condition-matched experimental benchmark. The modest population-level correlations and the sensitivity to baseline correction show that acquisition conditions and preprocessing remain consequential. The RRUFF development split also overlaps in identifier between optimization and checkpoint selection, and its external test is disjoint by identifier rather than mineral species. These results motivate prospective paired measurements and stricter validation partitions before experimental use.

Several methodological limitations are equally relevant. Sample-wise normalization defines a spectral-shape prediction task; reporting errors after rescaling by target extrema does not demonstrate recovery of unknown absolute intensities. Deterministic integration supplies a point estimate without calibrated uncertainty, particularly problematic when a source spectrum is compatible with multiple target patterns. Molecular overlap across data resources and the provenance of descriptor labels also require explicit accounting when assessing transfer. Future work should prioritize paired experimental data, calibration and uncertainty assessment, together with comparisons that match architectures, optimization budgets and inference cost.

In summary, the reported results support cross-modal spectral translation as a useful objective for normalized spectral completion and molecular feature learning. The next step is to establish how reliably these capabilities extend to calibrated measurements and unfamiliar molecular systems, while resolving the peak-level errors and protocol limitations identified here.
